
% ------------------------------------------------------------
\section{Numerical methods}
\label{sec:numerics}
% ------------------------------------------------------------

Throughout this section we fix $D_0=\kappa=1$ for concreteness; the dependence
of every quantity on these parameters is explicit in the analytical formulas.
The full Python/NumPy implementation is provided in the accompanying notebook
\texttt{mfpt\_finiteDifference\_markov.ipynb}.  We compare three independent
numerical methods against the analytical formula \eqref{eq:T-avg-tent}:

\begin{itemize}
\item A direct finite-difference solution of the backward equation
  \eqref{eq:backwardODE}.
\item A Crank--Nicolson finite-difference solution of the forward Fokker--Planck
  equation \eqref{eq:forward}, from which the MFPT is recovered via survival
  probabilities.
\item A Monte Carlo simulation of the It\^o SDE \eqref{eq:ito-sde}.
\end{itemize}

\subsection{Finite-difference solution of the backward equation}

The spatial grid is uniform with $N$ cells on the ring,
$x_j=-\pi+j\,h$, $h=2\pi/N$, $j=0,\dots,N-1$, identifying $x_N\equiv x_0$.
Midpoint diffusivities are $D_{j+1/2}=\tfrac12(D(x_j)+D(x_{j+1}))$.  The
conservative (flux-form) second-order discretisation of
$(D(x)T')'$ at node $j$ is
\begin{equation}
\mathcal{L}_h T
=\frac{D_{j+1/2}(T_{j+1}-T_j)-D_{j-1/2}(T_j-T_{j-1})}{h^2},
\label{eq:fd-stencil}
\end{equation}
which is a symmetric tridiagonal (with periodic wrap-around) and preserves
self-adjointness of $\mathcal{L}$.  The $\delta$-sink at $x=0$ is approximated
at the grid node $j_0$ nearest to $0$ by $-\kappa\,\delta(x)\,p\to-(\kappa/h)T_{j_0}$,
which adds $-\kappa/h$ to the diagonal entry $A_{j_0,j_0}$.  The resulting
linear system $A\,\mathbf{T}=-\mathbf{1}$ is solved with
\texttt{scipy.sparse.linalg.spsolve}.  Without the sink, $A$ would have a
one-dimensional null space (the constant vector); the sink breaks this kernel
and renders the system non-singular.

\subsection{Crank--Nicolson solution of the forward FPE}

The same discrete generator is used in the forward problem: with the
time-dependent state vector $\mathbf{p}(t)\in\mathbb{R}^N$, the semi-discrete
system is
\[
\frac{d\mathbf{p}}{dt} \;=\; \mathcal{L}_h\,\mathbf{p},
\]
where $\mathcal{L}_h$ includes the discrete $\delta$-sink at $j_0$.  Time
stepping is by Crank--Nicolson with step $\Delta t$:
\[
\Bigl(I - \tfrac{\Delta t}{2}\,\mathcal{L}_h\Bigr)\mathbf{p}^{n+1}
=\Bigl(I + \tfrac{\Delta t}{2}\,\mathcal{L}_h\Bigr)\mathbf{p}^{n}.
\]
The implicit solve uses a pre-factored sparse LU decomposition
(\texttt{scipy.sparse.linalg.factorized}); solving for multiple initial
conditions simultaneously amounts to applying the solver to the columns of a
single right-hand-side matrix.  The MFPT starting from $x_0$ is recovered from
the survival probability $S(x_0,t)=h\sum_j p_j(t\mid x_0)$ via
\[
T(x_0)=\int_0^\infty\!S(x_0,t)\,dt
\;\approx\; \operatorname{trapz}(S,t) + \text{tail correction}.
\]
The tail correction fits $\log S$ linearly on the last 20\% of the time grid
to extract the leading exponential decay rate and then integrates analytically
beyond $T_{\max}$.

\subsection{Monte Carlo simulation of the It\^o SDE}

Euler--Maruyama applied to \eqref{eq:ito-sde} gives
\[
x_{n+1} \;=\; x_n \;+\; D'(x_n)\,\Delta t \;+\; \sqrt{2\,D(x_n)\,\Delta t}\;\xi_n,
\qquad \xi_n\sim\mathcal{N}(0,1),
\]
with $D'(x_n)=-D_1\,\mathrm{sign}(x_n)/\pi$ and periodic wrapping of positions.
Partial absorption at the trap is implemented as a pointwise Poisson rule:
whenever $|x_n|<\varepsilon$, a particle is absorbed with probability
$p_{\rm abs}=\kappa\,\Delta t/(2\varepsilon)$ per step.  In the continuous
limit $(\Delta t,\varepsilon)\to 0$ with
$\varepsilon\gg\sqrt{2 D_{\max}\Delta t}$, the rule converges to the
continuous $\kappa\,\delta(x)$ sink.  We run $M$ particles with uniformly
random initial positions on $[-\pi,\pi]$ and record the first absorption
time $\tau^{(m)}$.  The empirical mean $\hat T = M^{-1}\sum_m \tau^{(m)}$
estimates $\langle T\rangle$; we also bin first-passage times by starting
position to obtain a position-resolved $T(x_0)$.
